\section{Problem Formulation}
\label{sec:problem}
We study \emph{state-channel toxicity propagation} in multi-agent LLM systems: how a single toxic injection persists through the evolving interaction state, what channels carry it, and what classes of mitigation can close them.
This formulation departs from per-message toxicity evaluation, which treats each generation as an isolated output, and from classical machine unlearning, which removes specific memorized examples from model parameters.
Both miss the central object of interest here: a stateful behavioral pattern in which agents absorb, propagate, and amplify toxic influence through the conversation graph, the compressed memory, and the model's own register-matching tendencies.

\subsection{Agent-Mediated Conversations as State Machines}
\label{subsec:setup_notation}
We model an agent-mediated discussion as a directed graph $G = (V, E)$ generated over a shared, evolving state.
The process begins with a human seed post $s$ and proceeds through a sequence of agent-generated messages.
Each node $v \in V$ represents a message $x_v$ produced by an agent $a_v \in \mathcal{A}$, and each edge $(u \to v) \in E$ indicates that $x_v$ was generated conditioned on $x_u$ .
The authorship map $v \mapsto a_v$ is \emph{many-to-one}: a single agent may author multiple nodes in the graph.
This matters for the focal agent $A_1$ defined below, which under the \emph{multi-injection} condition authors several replies along a thread rather than only the first response.

At each generation step $t$, the acting agent observes a \emph{conditioning set} $\mathcal{C}(v)$ drawn from the current state and produces a message:
\begin{equation}
    x_v \sim \pi_\theta(\cdot \mid \mathcal{C}(v), r_v),
\end{equation}
where $\pi_\theta$ is the agent's policy (parameterized by an LLM with parameters $\theta$), $\mathcal{C}(v)$ is the visible context, and $r_v$ denotes the agent's role.
The state is then updated to include $x_v$, making it available to future agents. This observe--generate--update loop creates a closed-loop dynamical system in which each agent's output becomes part of the conditioning context for subsequent agents.

\subsection{Toxic Influence as Stateful Contamination}
 
We define the \emph{focal agent} $A_1$ as the first responder to the seed post and vary $A_1$'s behavior:
\[
A_1 \in \{\texttt{neutral},\ \texttt{toxic}\},
\]
while keeping all other components fixed
% ---the seed post, the downstream agents $A_2, \ldots, A_n$, the graph topology, the generation schedule, and the decoding randomness.
, and each message is scored by a toxicity model (Detoxify~\citep{Detoxify}):
\begin{equation}
    \mathrm{tox}(v) = f_{\mathrm{tox}}(x_v) \in [0, 1].
\end{equation}
where larger values indicate more toxic content.
A standard filtering threshold $\tau{=}0.5$ flags individual messages as overtly toxic.

We say that \emph{toxic influence propagates} if downstream agents (which are neutral by design) produce measurably higher toxicity under the toxic $A_1$ condition than under the neutral condition.
To isolate downstream propagation from $A_1$'s own contribution, we restrict the average to nodes \emph{not authored by $A_1$}.
Let $V_{A_1} = \{v \in V : a_v = A_1\}$ denote the set of focal-agent nodes (a singleton in the single-injection case, larger under multi-injection).
The downstream toxicity of a graph is then
\begin{equation}
    \mu(G) = \frac{1}{|V \setminus V_{A_1}|} \sum_{v \in V \setminus V_{A_1}} \mathrm{tox}(v),
\end{equation}
and the \textbf{paired effect size} for seed $s$ is:
\begin{equation}
    \Delta\mu(s) = \mu\!\left(G^{(\texttt{toxic})}(s)\right) - \mu\!\left(G^{(\texttt{neutral})}(s)\right).
    \label{eq:effect_size}
\end{equation}
This definition guarantees that any positive $\Delta\mu$ reflects a change in non-focal agents' behavior, not the mechanically high toxicity of $A_1$'s own message.
A positive $\Delta\mu$ therefore indicates that $A_1$'s toxicity has causally influenced the downstream agents' generations.

\subsection{Three Channels of Toxic Persistence}
\label{sec:channels}
We identify three distinct channels through which toxic influence persists in agentic systems:

\noindent \textbf{(i) Transcript backflow.}
Toxic content produced by $A_1$ remains in the raw conversation history and re-enters downstream agents' context at every subsequent turn. Even if a downstream agent would not independently produce toxic content, conditioning on a transcript that contains a high-toxicity message can elevate the Detoxify score of its own output. This is the dominant channel in systems where agents see the full conversation history.

\noindent \textbf{(ii) Memory laundering.}
When agents maintain compressed memory representations (e.g., running summaries of the conversation), toxic framing can be \emph{laundered} through summarization.
The resulting summary may not score as explicitly toxic on classifiers, for instance, a summary might read ``\emph{the discussion has become heated, with participants expressing strong disagreement}'' --- yet conditioning on it raises the expected Detoxify score of subsequent generations relative to conditioning on the matched neutral summary.
We formalize this as:
\begin{equation}
    \mathrm{tox}(M_t) < \tau \quad \text{but} \quad \mathbb{E}\!\left[\mathrm{tox}(v) \mid M_t \in \mathcal{C}(v)\right] > \mathbb{E}\!\left[\mathrm{tox}(v) \mid M_t^{(\texttt{neutral})} \in \mathcal{C}(v)\right],
    \label{eq:laundered}
\end{equation}
where $M_t$ is the contaminated memory and $M_t^{(\texttt{neutral})}$ is the memory from the paired neutral run.

\noindent \textbf{(iii) Parametric bias.}
The model's own learned tendencies can interact with toxic context to amplify propagation.
LLMs tend to match the lexical and stylistic register of their input; when conditioned on a high-Detoxify-score context, this tendency yields outputs with higher Detoxify scores than the model would produce from a neutral context, even after safety fine-tuning.
This channel compounds with both \emph{transcript backflow} (channel i) and \emph{memory laundering} (channel ii): toxic content surfaced through either channel re-enters the model's input and triggers register-matching at the next generation step.

A robust mitigation must therefore prevent relapse under all three channels simultaneously --- a property we formalize as \emph{backflow-robust mitigation} in~\cref{app:formal_definitions}. Standard LLM unlearning methods (parameter editing, gradient ascent on forget sets) achieve $\Delta\mu \approx 0$ under single-turn evaluation but, as standalone mitigations, are not backflow-robust: they intervene only on \emph{parametric bias} (channel iii) and leave \emph{transcript backflow} (channel i) and \emph{memory laundering} (channel ii) open, allowing toxic content to re-enter the state and trigger relapse.

\subsection{Objective}
\label{sec:objective}
Our goal is to learn an updated system $(\pi_{\theta'}, S, W)$, where $\pi_{\theta'}$ is a fine-tuned policy that reduces parametric amplification of toxic context (e.g., via DPO), $S$ is a \textbf{read-side sanitizer} applied to the conditioning set $\mathcal{C}(v)$ before generation (redacts or rewrites toxic content surfaced through transcript or memory), and $W$ is a \textbf{write-side gate} applied after generation that controls whether the produced message and its memory update enter the state.
The objective is:
\begin{equation}
    \mathbb{E}_s\!\left[\Delta\mu^{(\text{post})}(s)\right] \approx 0
    \quad \text{subject to} \quad
    U(\tau') \approx U(\tau),
    \label{eq:objective}
\end{equation}
where $U(\cdot)$ measures conversational coherence, topical relevance, and response quality, ensuring that the mitigation does not degrade the agent's utility as a conversational participant.
 
The central challenge is that no single component suffices, and the failure mode of each maps onto a different channel.
Parameter-level unlearning ($\pi_{\theta'}$ alone) addresses \emph{parametric bias} (channel iii) but leaves transcript backflow and memory laundering open; toxic content re-enters the state and overwhelms the parametric correction.
State-level control ($S, W$ alone) closes channels (i) and (ii) but leaves channel (iii) open: residual toxic context that escapes sanitization can still trigger register-matching at generation.
Memory-only intervention addresses only channel (ii).
This motivates our \emph{three-pathway} framework, in which all three channels are addressed simultaneously.
 
% \paragraph{Section summary.}
% We have formalized agent-mediated discussion as a closed-loop state machine and identified three distinct channels --- \emph{transcript backflow}, \emph{memory laundering}, and \emph{parametric bias} --- through which a single toxic injection persists in the state and shapes downstream behavior.
% We defined the paired effect size $\Delta\mu$ over non-focal nodes as the propagation metric and \emph{backflow-robust mitigation} as the property that no robust defense can omit any one of the three channels.
% The next section instantiates this framework as a three-pathway intervention $(\pi_{\theta'}, S, W)$ and details how each component targets one channel.

% \subsection{Definitions}
 
% We formalize two key properties that a robust state-channel mitigation method must address:
 
% \begin{definition}[Relapse under re-exposure]
% \label{def:relapse}
% An agent exhibits \emph{relapse} if, after a mitigation intervention $\mathcal{I}$, it produces elevated toxicity when re-exposed to toxic context through the evolving state:
% \begin{equation}
%     \exists\, t > t_{\mathcal{I}}:\quad \mathrm{tox}(v_t) > \mathrm{tox}(v_t^{(\texttt{clean})}),
% \end{equation}
% where $v_t^{(\texttt{clean})}$ is the message the agent would produce under clean (intervention-maintained) state, and $v_t$ is the message produced when the state has been re-contaminated by upstream toxic content.
% \end{definition}
 
% \begin{definition}[Backflow-robust mitigation]
% \label{def:backflow_robust}
% A mitigation method is \emph{backflow-robust} if it prevents relapse under continuous re-exposure: for all $t$ after intervention,
% \begin{equation}
%     \Delta\mu^{(\text{post-}\mathcal{I})}(s) \approx 0,
% \end{equation}
% even when the environment continues to surface toxic context through transcript accumulation and memory updates.
% \end{definition}

% Standard LLM unlearning methods (parameter editing, gradient ascent on forget sets) achieve $\Delta\mu \approx 0$ under single-turn evaluation but, as standalone mitigations, fail to satisfy Definition~\ref{def:backflow_robust}: they intervene only on \emph{parametric bias} (channel iii) and leave \emph{transcript backflow} (channel i) and \emph{memory laundering} (channel ii) open, allowing toxic content to re-enter the state and trigger relapse.

